% ============================================================================
% COURS 04 — CINÉMATIQUE — PARTIE C
% Source : 5-Cinematique2022.docx
% ============================================================================

\section{Composition des mouvements}
\label{sec:cine-composition}

\subsection{Composition des vecteurs vitesse}

Pour un point $P$ appartenant au solide 2, lui-même en mouvement par
rapport au solide 1, en mouvement par rapport à 0 :
\[
\boxed{
\vec V(P\in2/0)
=\vec V(P\in2/1)+\vec V(P\in1/0).}
\]
Le premier terme est la \textbf{vitesse relative}; le second est la
\textbf{vitesse d'entraînement}, c'est-à-dire la vitesse du point $P$
considéré comme fixé au solide 1.

Pour $n$ solides successifs :
\[
\vec V(P\in n/0)
=\vec V(P\in n/n-1)+\cdots+\vec V(P\in1/0).
\]
Cette écriture est valable au même point géométrique $P$.

\subsection{Composition des torseurs cinématiques}

\[
\boxed{
\left\{\mathcal V_{n/0}\right\}_P
=\left\{\mathcal V_{n/n-1}\right\}_P+
\cdots+
\left\{\mathcal V_{1/0}\right\}_P.}
\]
Tous les torseurs doivent impérativement être exprimés au même point et
dans une base commune avant addition.

\begin{TLCineMethod}{Déterminer une vitesse par composition}
\begin{enumerate}
  \item Choisir une chaîne de solides reliant le solide étudié au bâti.
  \item Écrire la composition au point demandé.
  \item Pour chaque mouvement simple, déplacer le torseur vers ce point.
  \item Additionner les vecteurs rotation et les vecteurs vitesse.
  \item Laisser si possible le résultat dans les bases naturelles des mouvements.
  \item Vérifier les unités et la direction sur une position particulière.
\end{enumerate}
\end{TLCineMethod}

\subsection{Champ des vecteurs vitesse}

Pour deux points $A$ et $B$ d'un même solide 1 :
\[
\boxed{
\vec V(B\in1/0)
=\vec V(A\in1/0)+\overrightarrow{BA}\wedge\vec\Omega_{1/0}.}
\]
Cette relation ne doit jamais être utilisée entre deux points appartenant
à des solides différents.

\begin{TLCineApp}{Rotation autour d'un axe}
Si $A$ appartient à l'axe de rotation, alors
$\vec V(A\in1/0)=\vec0$ et :
\[
\vec V(B\in1/0)=\overrightarrow{BA}\wedge\vec\Omega_{1/0}.
\]
Le vecteur est tangent au cercle décrit par $B$, et sa norme est
$AB\,\Omega$ lorsque $AB$ est perpendiculaire à l'axe.
\end{TLCineApp}

\subsection{Accélération d'un point d'un solide}

La méthode recommandée consiste à dériver le vecteur vitesse dans le
repère de référence. Pour deux points d'un même solide :
\[
\vec A(B\in1/0)=\vec A(A\in1/0)
+\overrightarrow{BA}\wedge\dot{\vec\Omega}_{1/0}
+(\vec\Omega_{1/0}\wedge\overrightarrow{BA})
\wedge\vec\Omega_{1/0}.
\]
Le deuxième terme est tangentiel et le troisième centripète.

Pour un point appartenant à un solide mobile 2 dans un repère mobile 1 :
\[
\boxed{
\vec A(P\in2/0)=\vec A(P\in2/1)+\vec A(P\in1/0)
+2\vec\Omega_{1/0}\wedge\vec V(P\in2/1).}
\]
Le dernier terme est l'accélération de Coriolis.

\section{Fermeture cinématique}
\label{sec:cine-fermeture}

Dans une boucle cinématique fermée :
\[
\boxed{
\left\{\mathcal V_{1/0}\right\}
+\left\{\mathcal V_{2/1}\right\}
+\cdots+
\left\{\mathcal V_{0/n}\right\}
=\left\{0\right\}.}
\]
Après réduction au même point, les six composantes fournissent un
système d'équations. Cette fermeture permet d'obtenir une loi
entrée-sortie en vitesse, d'étudier les mobilités et d'évaluer le degré
d'hyperstatisme.

\begin{TLCineApp}{Vérin électrique}
Une vis et un écrou sont liés par une liaison hélicoïdale et guidés par
une pivot et une glissière par rapport au bâti. Écrire la fermeture :
\[
\left\{\mathcal V_{2/0}\right\}
-\left\{\mathcal V_{1/0}\right\}
-\left\{\mathcal V_{2/1}\right\}=\{0\}.
\]
Projeter les composantes utiles pour retrouver
$v=(p/2\pi)\omega$ et identifier les équations supplémentaires
traduisant les contraintes internes.
\end{TLCineApp}

\section{Cinématique du contact ponctuel}
\label{sec:cine-contact}

Lorsque deux solides 1 et 2 sont en contact au point géométrique $I$, il
faut distinguer :
\begin{itemize}
  \item le point matériel $I_1$ appartenant au solide 1 ;
  \item le point matériel $I_2$ appartenant au solide 2 ;
  \item le point géométrique de contact $I$.
\end{itemize}

La vitesse de glissement de 2 par rapport à 1 est :
\[
\boxed{
\vec V(I\in2/1)
=\vec V(I\in2/0)-\vec V(I\in1/0).}
\]
Elle appartient au plan tangent commun aux deux solides.

\begin{TLCineDef}{Roulement sans glissement}
Il y a roulement sans glissement au point $I$ lorsque :
\[
\boxed{\vec V(I\in2/1)=\vec0.}
\]
Le point $I$ est alors un centre instantané de rotation pour un
mouvement plan. Cette condition est utilisée pour les engrenages,
roulements, roues sur le sol, poulies-courroies et pignons-chaînes.
\end{TLCineDef}

\begin{TLCineApp}{Roue sur un plateau}
Une roue 2 entraînée par rapport au bras 3 roule sur un plateau 1,
lui-même en rotation par rapport au bâti 0. Écrire :
\[
\vec V(I\in2/1)=\vec V(I\in2/0)-\vec V(I\in1/0)=\vec0.
\]
Décomposer ensuite chacune des vitesses par composition pour relier les
trois vitesses de rotation.
\end{TLCineApp}

\section{Loi de vitesse en trapèze}
\label{sec:cine-trapeze}

Un actionneur ne peut pas passer instantanément d'une vitesse nulle à
une vitesse non nulle. Une loi simple de commande comporte :
\begin{enumerate}
  \item une phase d'accélération constante ;
  \item une phase à vitesse constante ;
  \item une phase de décélération constante.
\end{enumerate}

\begin{center}
\begin{tikzpicture}[>=Latex,x=1.35cm,y=1.2cm]
  \draw[->] (0,0)--(7.2,0) node[right]{$t$};
  \draw[->] (0,0)--(0,3.2) node[above]{$v(t)$};
  \draw[very thick,TLCineBlue]
    (0,0)--(1.6,2.5)--(4.8,2.5)--(6.5,0);
  \draw[dashed] (1.6,0)--(1.6,2.5);
  \draw[dashed] (4.8,0)--(4.8,2.5);
  \draw[dashed] (6.5,0)--(6.5,2.5);
  \node[below] at (1.6,0) {$t_1$};
  \node[below] at (4.8,0) {$t_2$};
  \node[below] at (6.5,0) {$t_3$};
  \node[left] at (0,2.5) {$V$};
\end{tikzpicture}
\end{center}

Les relations fondamentales sont :
\[
v(t)=\dot x(t),
\qquad
a(t)=\dot v(t)=\ddot x(t),
\]
\[
\boxed{x(t)-x(0)=\int_0^t v(T)\,\mathrm dT.}
\]
Le déplacement est donc l'aire algébrique sous la courbe de vitesse.
L'accélération est la pente de cette courbe.

Pour une vitesse maximale $V$, avec accélération de durée $t_a$, palier
de durée $t_p$ et décélération de durée $t_d$, le déplacement total est :
\[
\boxed{
\Delta x=\frac12Vt_a+Vt_p+\frac12Vt_d.}
\]
Si $t_a=t_d$ :
\[
\Delta x=V(t_p+t_a).
\]

\begin{TLCineMethod}{Exploiter un profil de vitesse}
\begin{enumerate}
  \item Lire les instants de changement de phase et les vitesses.
  \item Calculer les accélérations avec $a=\Delta v/\Delta t$.
  \item Déterminer les déplacements par les aires : triangles, rectangles et trapèzes.
  \item Construire la courbe de position, continue et de pente égale à la vitesse.
  \item Vérifier la continuité de la vitesse et de la position aux changements de phase.
\end{enumerate}
\end{TLCineMethod}

\begin{TLCineApp}{Hayon électrique}
Le hayon atteint $20^\circ\,\mathrm{s^{-1}}$ en $0{,}5\ \mathrm s$,
conserve cette vitesse puis décélère symétriquement. Déterminer la durée
du palier pour une ouverture totale de $90^\circ$, puis tracer
$\gamma(t)$ et $\ddot\gamma(t)$.
\end{TLCineApp}

\section{Sources}
\label{sec:cine-sources}

Ce cours reprend le document original, les systèmes pédagogiques
associés et des ressources de collègues de l'UPSTI. Les exercices sont
issus de sujets de concours ou d'activités explicitement mentionnés dans
le document source.
